\documentclass[../main-en.tex]{subfiles}
\begin{document}
\renewcommand{\baselinestretch}{1.5}
\section{Architectural Design Rationale}
\label{sec:rationale}

% Format: name the decision, the tactic(s) used, and why — "this tactic is
% used because...". Keep each rationale short and tied to a quality attribute.

\begin{table}[h]
  \centering
  \begin{tabularx}{\textwidth}{| X | X | X |}
    \hline
    \textbf{Decision} & \textbf{Tactic} & \textbf{Rationale} \\
    \hline
    Sources, transformations and sinks sit behind interfaces & Encapsulate &
    This tactic is used because new sources and sinks must be added without
    modifying the pipeline; changes stay local to one module. \\
    \hline
    The transformation chain (incl. unit conversion) is defined in a config
    file & Use configuration parameters, defer binding &
    This tactic is used because the chain must be changed at deployment
    time without recompiling the code. \\
    \hline
    % TODO: add remaining key decisions, e.g. how sinks receive results,
    % how readings flow through the pipeline, choice of data format.
    \hline
  \end{tabularx}
  \caption{Key architectural decisions with their tactics and rationale.}
  \label{tab:rationale}
\end{table}

\end{document}
